\section{Amplitude equations and patterns}

Considering only non-homogeneous forcing at $2\omega_0$ and $2k_0$, define
the complex forcing coefficients
\[
 \eta_t\coloneqq
 2\Psi^*\cQ\!\left(\overline\Phi,p_{2\omega_0}\right),
 \qquad
 \eta_x\coloneqq
 2\Psi^*\cQ\!\left(\Phi,s_{2k_0}\right).
\]
Thus $\eta_t=0$ when $P=0$, and $\eta_x=0$ when $S=0$.  The amplitude
equations are
\begin{subequations}
\label{eq:amplitude}
\begin{align}
 A_\tau&=A\left(a+b|A|^2+c|B|^2\right)
          +\eta_t\overline B+\eta_xB,
 \label{eq:A}\\
 B_\tau&=B\left(a+b|B|^2+c|A|^2\right)
          +\eta_t\overline A+\eta_xA.
 \label{eq:B}
\end{align}
\end{subequations}
Here $a,b,c,\eta_t,\eta_x\in\mathbb C$.  An equilibrium satisfies
\begin{align}
 0&=A\left(a+b|A|^2+c|B|^2\right)
       +\eta_t\overline B+\eta_xB,
 \label{eq:A-equilibrium}\\
 0&=B\left(a+b|B|^2+c|A|^2\right)
       +\eta_t\overline A+\eta_xA.
 \label{eq:B-equilibrium}
\end{align}

\subsection{Nonzero equilibrium solution}

Assume first that $A,B\neq0$.  Set
\[
 x=|A|^2,
 \qquad
 \zeta=\frac BA,
 \qquad
 u=\frac{\overline A}{A}.
\]
Then $|u|=1$, $|B|^2=|\zeta|^2x$, and, for some $\phi\in\mathbb R$,
\[
 A=\sqrt{x}\,\ee^{\ii\phi},
 \qquad
 u=\ee^{-2\ii\phi},
 \qquad
 B=\zeta A.
\]
Dividing \eqref{eq:A-equilibrium} by $A$, and dividing
\eqref{eq:B-equilibrium} by $A$, gives
\begin{align}
 0&=a+bx+c|\zeta|^2x
      +\eta_t\overline\zeta u+\eta_x\zeta,
 \label{eq:ratio-A}\\
 0&=\zeta\left(a+b|\zeta|^2x+cx\right)
      +\eta_tu+\eta_x.
 \label{eq:ratio-B}
\end{align}
Eliminating $u$ gives
\begin{equation}
 \left(1-|\zeta|^2\right)
 \left[a+bx\left(1+|\zeta|^2\right)\right]
 +\eta_x\left(\zeta-\overline\zeta\right)=0.
 \label{eq:ratio-balance}
\end{equation}
The condition $|u|=1$ gives
\begin{equation}
 \left|
 \eta_x+\zeta\left[a+\left(b|\zeta|^2+c\right)x\right]
 \right|=|\eta_t|.
 \label{eq:ratio-modulus}
\end{equation}
Equations \eqref{eq:ratio-balance} and \eqref{eq:ratio-modulus} are
necessary and sufficient.  The following reduction gives their closed-form
solution.

Define the real invariants
\[
 x=|A|^2,
 \qquad
 T=x+|B|^2,
 \qquad
 D=x-|B|^2,
\]
and
\[
 A\overline B=h+\ii k,
 \qquad
 h,k\in\mathbb R.
\]
They satisfy
\begin{equation}
 h^2+k^2=x(T-x).
 \label{eq:relative-phase-constraint}
\end{equation}
Multiplying \eqref{eq:A-equilibrium} by $\overline A$, multiplying
\eqref{eq:B-equilibrium} by $\overline B$, and subtracting gives
\begin{equation}
 D(a+bT)-2\ii\eta_xk=0.
 \label{eq:difference-balance}
\end{equation}

\subsubsection{Unequal intensities}

Suppose $D\neq0$, hence $|A|\neq|B|$.  We consider the three forcing cases
separately.

\paragraph{Only spatial forcing: $\eta_t=0$ and $\eta_x\neq0$.}
Set
\[
 K=\frac{k}{D}\in\mathbb R.
\]
Equation \eqref{eq:difference-balance} gives
\begin{equation}
 a+bT-2\ii\eta_xK=0,
 \label{eq:spatial-TK-equation}
\end{equation}
and hence
\begin{equation}
 T=-\frac{\operatorname{Re}(\overline{\eta_x}a)}
          {\operatorname{Re}(\overline{\eta_x}b)},
 \qquad
 K=\frac{\operatorname{Im}(\overline ba)}
          {2\operatorname{Re}(\overline b\eta_x)}.
 \label{eq:spatial-TK-solution}
\end{equation}
The remaining equations are
\begin{align}
 h^2
 &=\frac{T^2}{4}-\left(K^2+\frac14\right)D^2,
 \label{eq:spatial-h-constraint}\\
 0
 &=\frac{|c-b|^2}{4}(T^2-D^2)
   +2\operatorname{Re}\!\left(\overline{\eta_x}(c-b)\right)h
   +(1+4K^2)|\eta_x|^2
   +2TK\operatorname{Im}\!\left(\overline{\eta_x}(c-b)\right).
 \label{eq:spatial-h-balance}
\end{align}
If $\operatorname{Re}(\overline{\eta_x}(c-b))\neq0$, eliminate $h$ to
obtain
\begin{equation}
 \begin{aligned}
 &\left[
  \begin{aligned}
   &\frac{|c-b|^2}{4}(T^2-D^2)+(1+4K^2)|\eta_x|^2\\
   &\qquad
    +2TK\operatorname{Im}\!\left(\overline{\eta_x}(c-b)\right)
  \end{aligned}
 \right]^2\\
 &\qquad
 =4\operatorname{Re}\!\left(\overline{\eta_x}(c-b)\right)^2
 \left[
  \frac{T^2}{4}-\left(K^2+\frac14\right)D^2
 \right].
 \end{aligned}
 \label{eq:spatial-D-quadratic}
\end{equation}
This is quadratic in $D^2$.  Retain only real roots satisfying
\begin{equation}
 0<D^2<T^2,
 \qquad
 \frac{T^2}{4}-\left(K^2+\frac14\right)D^2\geq0.
 \label{eq:spatial-admissibility}
\end{equation}
For each root, take both signs of $D$ and compute
\begin{equation}
 \begin{aligned}
 x&=\frac{T+D}{2},
 &k&=KD,\\
 h&=-\frac{
 \dfrac{|c-b|^2}{4}(T^2-D^2)+(1+4K^2)|\eta_x|^2
 +2TK\operatorname{Im}\!\left(\overline{\eta_x}(c-b)\right)}
 {2\operatorname{Re}\!\left(\overline{\eta_x}(c-b)\right)},
 &\zeta&=\frac{h-\ii k}{x}.
 \end{aligned}
 \label{eq:spatial-reconstruction}
\end{equation}
Then
\begin{equation}
 \eta_x+\zeta\left[a+b(T-x)+cx\right]=0,
 \label{eq:spatial-ratio}
\end{equation}
and the equilibria are
\begin{equation}
 A=\sqrt{x}\,\ee^{\ii\phi},
 \qquad
 B=\zeta A,
 \qquad
 \phi\in\mathbb R.
 \label{eq:spatial-equilibria}
\end{equation}
The arbitrary common phase follows from the spatial-forcing phase symmetry.

\paragraph{Only temporal forcing: $\eta_t\neq0$ and $\eta_x=0$.}
Equation \eqref{eq:difference-balance} gives
\begin{equation}
 T=-\frac ab\in\mathbb R_{>0}.
 \label{eq:temporal-T}
\end{equation}
The modulus condition reduces to
\begin{equation}
 |c-b|^2x(T-x)=|\eta_t|^2.
 \label{eq:temporal-x-balance}
\end{equation}
For $c\neq b$, the two exchange-related intensities are
\begin{equation}
 x_\pm=
 \frac{T\pm\sqrt{T^2-4|\eta_t|^2/|c-b|^2}}{2}.
 \label{eq:temporal-x-roots}
\end{equation}
Unequal-intensity equilibria exist when
\begin{equation}
 T^2>\frac{4|\eta_t|^2}{|c-b|^2}.
 \label{eq:temporal-existence}
\end{equation}
For each $x=x_\pm$, choose an arbitrary $\delta\in\mathbb R$ and set
\begin{equation}
 \zeta=\sqrt{\frac{T-x}{x}}\,\ee^{\ii\delta},
 \qquad
 \ee^{-2\ii\phi}=-\frac{\zeta(c-b)x}{\eta_t}.
 \label{eq:temporal-phase}
\end{equation}
The equilibria are
\begin{equation}
 A=\sqrt{x}\,\ee^{\ii\phi},
 \qquad
 B=\zeta A.
 \label{eq:temporal-equilibria}
\end{equation}
The free value of $\delta$ represents the opposite-phase symmetry.

\paragraph{Combined forcing: $\eta_t\neq0$ and $\eta_x\neq0$.}
Compute $T$ and $K$ from \eqref{eq:spatial-TK-solution}.  The remaining
equations are
\begin{align}
 h^2
 &=\frac{T^2}{4}-\left(K^2+\frac14\right)D^2,
 \label{eq:combined-h-constraint}\\
 0
 &=\frac{|c-b|^2}{4}(T^2-D^2)
   +2\operatorname{Re}\!\left(\overline{\eta_x}(c-b)\right)h
   +|\eta_x|^2-|\eta_t|^2+4K^2|\eta_x|^2
   +2TK\operatorname{Im}\!\left(\overline{\eta_x}(c-b)\right).
 \label{eq:combined-h-balance}
\end{align}
If $\operatorname{Re}(\overline{\eta_x}(c-b))\neq0$, eliminate $h$ to
obtain
\begin{equation}
 \begin{aligned}
 &\left[
  \begin{aligned}
   &\frac{|c-b|^2}{4}(T^2-D^2)
    +|\eta_x|^2-|\eta_t|^2+4K^2|\eta_x|^2\\
   &\qquad
    +2TK\operatorname{Im}\!\left(\overline{\eta_x}(c-b)\right)
  \end{aligned}
 \right]^2\\
 &\qquad
 =4\operatorname{Re}\!\left(\overline{\eta_x}(c-b)\right)^2
 \left[
  \frac{T^2}{4}-\left(K^2+\frac14\right)D^2
 \right].
 \end{aligned}
 \label{eq:combined-D-quadratic}
\end{equation}
Retain only roots satisfying the conditions in
\eqref{eq:spatial-admissibility}.  For each root, take both signs of $D$ and
compute
\begin{equation}
 \begin{aligned}
 x&=\frac{T+D}{2},
 &k&=KD,\\
 h&=-\frac{
 \dfrac{|c-b|^2}{4}(T^2-D^2)
 +|\eta_x|^2-|\eta_t|^2+4K^2|\eta_x|^2
 +2TK\operatorname{Im}\!\left(\overline{\eta_x}(c-b)\right)}
 {2\operatorname{Re}\!\left(\overline{\eta_x}(c-b)\right)},
 &\zeta&=\frac{h-\ii k}{x}.
 \end{aligned}
 \label{eq:combined-reconstruction}
\end{equation}
The phase is fixed by
\begin{equation}
 \ee^{-2\ii\phi}
 =-\frac{\eta_x+\zeta\left[a+b(T-x)+cx\right]}{\eta_t},
 \label{eq:combined-phase}
\end{equation}
and the equilibria are
\begin{equation}
 A=\sqrt{x}\,\ee^{\ii\phi},
 \qquad
 B=\zeta A.
 \label{eq:combined-equilibria}
\end{equation}
Thus both forcing terms generically fix the two phases.

\paragraph{Consistency.}
The preceding formulas assume nonzero denominators and admissible real
roots.  If $\operatorname{Re}(\overline b\eta_x)=0$ in a case with
$\eta_x\neq0$, solve
\begin{equation}
 \begin{pmatrix}
  \operatorname{Re}b&2\operatorname{Im}\eta_x\\
  \operatorname{Im}b&-2\operatorname{Re}\eta_x
 \end{pmatrix}
 \begin{pmatrix}T\\K\end{pmatrix}
 =-
 \begin{pmatrix}\operatorname{Re}a\\\operatorname{Im}a\end{pmatrix}.
 \label{eq:TK-degenerate}
\end{equation}
If $\operatorname{Re}(\overline{\eta_x}(c-b))=0$, use the corresponding
unsquared balance equation to solve for $D^2$, and then use the constraint
equation to compute $h$.  An inconsistent system, $T\leq0$, or the absence
of an admissible root gives no unequal-intensity equilibrium.  A
rank-deficient consistent system may give a continuous family.  In the
temporal-only case, $a+bT=0$ must have a positive real solution.  If $b=0$
and $a\neq0$, no such equilibrium exists; if $a=b=0$, $T$ is not fixed by
this equation and a degenerate family may occur.  Also, $c=b$ gives no
unequal-intensity equilibrium when $\eta_t\neq0$.

\subsubsection{Equal intensities}

Suppose $D=0$, so $|A|=|B|=r$.  Equation
\eqref{eq:difference-balance} gives $\eta_xk=0$.

If $\eta_x\neq0$, then $k=0$ and hence
\begin{equation}
 B=\sigma A,
 \qquad
 \sigma\in\{1,-1\}.
 \label{eq:parity-state}
\end{equation}
Set
\[
 d=b+c,
 \qquad
 g_\sigma=a+\sigma\eta_x,
 \qquad
 y=r^2.
\]
The equilibrium condition is
\begin{equation}
 g_\sigma+dy+\sigma\eta_t\ee^{-2\ii\phi}=0.
 \label{eq:parity-balance}
\end{equation}
For $d\neq0$, define
\begin{equation}
 \Delta_\sigma=
 \operatorname{Re}(g_\sigma\overline d)^2
 -|d|^2\left(|g_\sigma|^2-|\eta_t|^2\right).
 \label{eq:parity-discriminant}
\end{equation}
The possible amplitudes are
\begin{equation}
 y_{\sigma,\pm}=
 \frac{-\operatorname{Re}(g_\sigma\overline d)
       \pm\sqrt{\Delta_\sigma}}{|d|^2}.
 \label{eq:parity-amplitudes}
\end{equation}
Retain only $\Delta_\sigma\geq0$ and $y_{\sigma,\pm}>0$.  If
$\eta_t\neq0$, the phase is determined by
\begin{equation}
 \ee^{-2\ii\phi}
 =-\frac{g_\sigma+dy}{\sigma\eta_t}.
 \label{eq:parity-phase}
\end{equation}
If $\eta_t=0$, require $g_\sigma+dy=0$, and $\phi$ is arbitrary.  If
$d=0$, use \eqref{eq:parity-balance} directly; it gives either no solution
or a continuum of amplitudes.

If $\eta_x=0$, write
\[
 A=r\ee^{\ii\phi},
 \qquad
 B=r\ee^{\ii\psi},
 \qquad
 \Theta=\phi+\psi.
\]
The equilibrium condition is
\begin{equation}
 a+(b+c)r^2+\eta_t\ee^{-\ii\Theta}=0.
 \label{eq:temporal-equal-balance}
\end{equation}
Thus $|a+(b+c)r^2|=|\eta_t|$.  For $b+c\neq0$, the possible amplitudes are
\begin{equation}
 r_\pm^2=
 \frac{-\operatorname{Re}\!\left(a\overline{(b+c)}\right)
 \pm\sqrt{
 \operatorname{Re}\!\left(a\overline{(b+c)}\right)^2
 -|b+c|^2\left(|a|^2-|\eta_t|^2\right)}}{|b+c|^2}.
 \label{eq:temporal-equal-amplitudes}
\end{equation}
For $\eta_t\neq0$, the phase sum is fixed by
\begin{equation}
 \ee^{-\ii\Theta}
 =-\frac{a+(b+c)r^2}{\eta_t},
 \label{eq:temporal-equal-phase}
\end{equation}
while $\phi-\psi$ is arbitrary.  If $\eta_t=0$, require
$a+(b+c)r^2=0$, and both phases are arbitrary.

\subsection{One-zero equilibrium solution}

Let $A=0$ and $B=r\ee^{\ii\psi}\neq0$.  The equilibrium equations reduce
to
\begin{equation}
 a+br^2=0,
 \qquad
 \eta_t\ee^{-2\ii\psi}+\eta_x=0.
 \label{eq:B-only-balance}
\end{equation}
Hence, for $\eta_t\neq0$, this branch exists precisely when
\begin{equation}
 r^2=-\frac ab\in\mathbb R_{>0},
 \qquad
 |\eta_t|=|\eta_x|,
 \qquad
 \ee^{-2\ii\psi}=-\frac{\eta_x}{\eta_t}.
 \label{eq:B-only-solution}
\end{equation}
If $\eta_t=0$, a one-zero equilibrium requires $\eta_x=0$ and
$a+br^2=0$; its active phase is arbitrary.  The solution with $B=0$ and
$A\neq0$ is obtained by exchanging $A$ and $B$.

\subsection{Zero equilibrium solution}

The zero equilibrium
\begin{equation}
 (A,B)=(0,0)
 \label{eq:zero-equilibrium}
\end{equation}
always exists.

\subsection{Stability}

Stability here refers to the amplitude equations in slow time $\tau$.  Write
$\mathbf A=(A,B)^\mathsf{T}$.  At an equilibrium $(A_*,B_*)$, the
linearization has
the form
\begin{equation}
 \delta\mathbf A_\tau
 =M\,\delta\mathbf A+N\,\delta\overline{\mathbf A},
 \label{eq:complex-linearization}
\end{equation}
where
\begin{equation}
 M=
 \begin{pmatrix}
  a+2b|A_*|^2+c|B_*|^2&cA_*\overline{B_*}+\eta_x\\
  cB_*\overline{A_*}+\eta_x&a+2b|B_*|^2+c|A_*|^2
 \end{pmatrix},
 \label{eq:M-block}
\end{equation}
and
\begin{equation}
 N=
 \begin{pmatrix}
  bA_*^2&cA_*B_*+\eta_t\\
  cA_*B_*+\eta_t&bB_*^2
 \end{pmatrix}.
 \label{eq:N-block}
\end{equation}
Since the system depends on complex conjugates, the eigenvalues of $M$ alone
do not determine stability.  The correct four-dimensional real
linearization is represented by
\begin{equation}
 \mathcal J=
 \begin{pmatrix}
  M&N\\
  \overline N&\overline M
 \end{pmatrix}.
 \label{eq:doubled-Jacobian}
\end{equation}
An isolated equilibrium is asymptotically stable if and only if every
eigenvalue of $\mathcal J$ has negative real part.

For a direct algebraic test, write
\begin{equation}
 \det(\lambda I-\mathcal J)
 =\lambda^4+\alpha_1\lambda^3+\alpha_2\lambda^2
  +\alpha_3\lambda+\alpha_4.
 \label{eq:quartic-characteristic}
\end{equation}
The coefficients $\alpha_j$ are real.  Strict stability is equivalent to
\begin{equation}
 \alpha_1,\alpha_2,\alpha_3,\alpha_4>0,
 \qquad
 \alpha_1\alpha_2>\alpha_3,
 \qquad
 \alpha_1\alpha_2\alpha_3>
 \alpha_3^2+\alpha_1^2\alpha_4.
 \label{eq:quartic-Hurwitz}
\end{equation}
This gives the stability of the unequal-intensity nonzero equilibria and the
one-zero equilibria after substituting their closed-form values into
\eqref{eq:M-block}--\eqref{eq:N-block}.

For the equal-intensity state $B_*=\sigma A_*$, perturbations split into
$\delta B=\kappa\delta A$, where $\kappa\in\{1,-1\}$.  Set
\begin{align}
 P_{\sigma\kappa}
 &=a+\kappa\eta_x+
   \left(2b+c+\kappa\sigma c\right)|A_*|^2,
 \label{eq:P-parity}\\
 Q_{\sigma\kappa}
 &=\left(b+\kappa\sigma c\right)A_*^2+\kappa\eta_t.
 \label{eq:Q-parity}
\end{align}
The two eigenvalues in this real two-dimensional block satisfy
\begin{equation}
 \lambda^2-2\operatorname{Re}(P_{\sigma\kappa})\lambda
 +|P_{\sigma\kappa}|^2-|Q_{\sigma\kappa}|^2=0.
 \label{eq:parity-characteristic}
\end{equation}
The equilibrium is stable if and only if, for both $\kappa=1$ and
$\kappa=-1$,
\begin{equation}
 \operatorname{Re}(P_{\sigma\kappa})<0,
 \qquad
 |P_{\sigma\kappa}|^2>|Q_{\sigma\kappa}|^2.
 \label{eq:parity-stability}
\end{equation}

For the one-zero equilibrium $(0,B_*)$, the blocks reduce to
\begin{equation}
 M_0=
 \begin{pmatrix}
  a+c|B_*|^2&\eta_x\\
  \eta_x&a+2b|B_*|^2
 \end{pmatrix},
 \qquad
 N_0=
 \begin{pmatrix}
  0&\eta_t\\
  \eta_t&bB_*^2
 \end{pmatrix}.
 \label{eq:one-zero-blocks}
\end{equation}
Its stability follows from \eqref{eq:doubled-Jacobian} and
\eqref{eq:quartic-Hurwitz}.  The case $(A_*,0)$ is obtained by exchanging the
two amplitudes.

At the zero equilibrium, define $Z_\sigma=A+\sigma B$ for
$\sigma\in\{1,-1\}$.  The two exchange-parity blocks are
\begin{equation}
 (Z_\sigma)_\tau
 =(a+\sigma\eta_x)Z_\sigma
  +\sigma\eta_t\overline{Z_\sigma}.
 \label{eq:zero-parity-block}
\end{equation}
Therefore the zero equilibrium is stable if and only if, for both
$\sigma=1$ and $\sigma=-1$,
\begin{equation}
 \operatorname{Re}(a+\sigma\eta_x)<0,
 \qquad
 |a+\sigma\eta_x|^2>|\eta_t|^2.
 \label{eq:zero-stability}
\end{equation}

For spatial forcing only, the common-phase symmetry gives one zero
eigenvalue; orbital stability requires the other three eigenvalues to have
negative real part.  For temporal forcing only, the opposite-phase symmetry
gives the same conclusion.  When both forcing coefficients are nonzero, both
phases are generically locked, and asymptotic stability requires all four
eigenvalues to have negative real part.  These amplitude stability results
give the leading local stability of the associated forced patterns near
onset.
